% Part: first-order-logic
% Chapter: model-theory
% Section: overspill

\documentclass[../../../include/open-logic-section]{subfiles}

\begin{document}

\olfileid[fr]{mod}{bas}{ove}
\olsection{Débordement}

\begin{thm}
\ollabel{overspill}
Si un ensemble $\Gamma$ d'énoncés possède des modèles finis de
cardinalités arbitrairement grandes, il possède un modèle infini.
\end{thm}

\begin{proof}
Élargissons le langage de $\Gamma$ en lui ajoutant une infinité
dénombrable de nouvelles constantes $c_0$, $c_1$, \dots, et
considérons l'ensemble $\Gamma \cup \{c_i \neq c_j : i \neq j\}$.
Dire que $\Gamma$ possède des modèles finis de cardinalités
arbitrairement grandes signifie que, pour tout $m >0$, il existe
$n\ge m$ tel que $\Gamma$ ait un modèle de cardinalité~$n$.
Il s'ensuit que tout sous-ensemble fini de $\Gamma \cup
\{c_i \neq c_j : i \neq j\}$ est satisfaisable. Par le théorème de
compacité, $\Gamma \cup \{c_i \neq c_j : i \neq j\}$ possède un
modèle $\Struct M$ dont le
!!{domain} est nécessairement infini, puisqu'il satisfait toutes
les inégalités $c_i \neq c_j$.
\end{proof}

\begin{prop}
\ollabel{inf-not-fo}
Il n'existe, dans aucun langage du premier ordre, d'énoncé $!A$ qui
soit vrai dans !!a{structure}~$\Struct M$ si et seulement si le
!!{domain} $\Domain{M}$ de cette !!{structure} est infini.
\end{prop}

\begin{proof}
S'il existait un tel énoncé~$!A$, sa négation $\lnot !A$ serait vraie
dans toutes les !!{structure}s finies et seulement dans celles-ci.
Elle aurait donc des modèles finis de cardinalités arbitrairement
grandes, mais aucun modèle infini, ce qui contredirait
\olref{overspill}.
\end{proof}

\end{document}
